\section{Screening for every external coupling}\label{app:D.27}

This section derives the screened response \cref{eq:8.4-screened} for any pair of target rows, its total-charge and electron-density forms, the exactness theorem \cref{eq:8.4-exact} in full, and the classification of Coulomb lines of \cref{sec:8.4-lines}. The notation is that of \cref{sec:4.4,sec:8.4}: $\vect{q}$ lies in the first zone, the reciprocal-lattice vectors $\vect{G}$ label the channel $\{\vect{q} + \vect{G}\}$, $v_{\vect{G}} \equiv v_{\mathrm{C}}(|\vect{q}+\vect{G}|)$, and the dark-matter vertex sits at the row $\vect{G}_q$ of the unreduced momentum (\cref{eq:A.3-unreduced}).

\subsection{The screened response}

The Coulomb Hamiltonian \cref{eq:2.3-HC} couples the species only through $n_Z = -n_e + \sum_\kappa Z_\kappa n_\kappa$, by the term $(2\Vol)^{-1}\sum_{\vect{k}\ne0}v_{\mathrm{C}}(k)\,n_Z(\vect{k})^\dagger n_Z(\vect{k})$ (the self-energies of the point charges aside). In the presence of a weak external field the induced charge-number density at $\vect{q} + \vect{G}$, $\delta\nd_{Z,\vect{G}} = \delta\langle n_Z(\vect{q}+\vect{G})^\dagger\rangle/\Vol$, adds to the Hamiltonian the mean-field term $v_{\vect{G}}\,\delta\nd_{Z,\vect{G}}\,n_Z(\vect{q}+\vect{G})$ plus its conjugate, which is a field of amplitude $v_{\vect{G}}\delta\nd_{Z,\vect{G}}$ on the row $n_Z(\vect{q}+\vect{G})$. The proper response of the channel is defined by \cref{eq:8.4-proper} as the response to the external fields $\phi_b$ and to these Hartree fields together. Everything else that the Coulomb interaction does, including the Coulomb lines at momenta outside the channel and the fluctuations of the channel components beyond their mean, stays inside $\tilde R$. By the argument of \cref{sec:4.8-removal}, $\tilde R$ is the response of the Hamiltonian with the Coulomb terms of the channel removed, in the thermodynamic limit.

Write the induced charge as a column vector $\delta\nd_Z$ over $\vect{G}$, the proper charge--charge block as the matrix $P_{\vect{G}\vect{G}'} = \tilde R_{Z_{\vect{G}}Z_{\vect{G}'}}$, and $v$ as the diagonal matrix of $v_{\vect{G}}$. For a field on row $b$ alone, \cref{eq:8.4-proper} with $a = Z_{\vect{G}}$ reads
\begin{equation}
  \delta\nd_Z = \tilde R_{Zb}\,\phi_b + P\,v\,\delta\nd_Z , \qquad
  \delta\nd_Z = (1 - Pv)^{-1}\tilde R_{Zb}\,\phi_b ,
  \label{eq:D.27-charge}
\end{equation}
with $\tilde R_{Zb}$ the column of $\tilde R_{Z_{\vect{G}}b}$. Inserting this into \cref{eq:8.4-proper} for a general row $a$,
\begin{equation*}
  \delta O_a = \bigl[\tilde R_{ab} + \tilde R_{aZ}\,v\,(1 - Pv)^{-1}\tilde R_{Zb}\bigr]\phi_b ,
\end{equation*}
with $\tilde R_{aZ}$ the row vector of $\tilde R_{aZ_{\vect{G}}}$. The push-through identity $v(1 - Pv)^{-1} = (1 - vP)^{-1}v$, which follows from $(1 - vP)v = v(1 - Pv)$ by multiplying with the two inverses, turns the bracket into
\begin{equation}
  R_{ab} = \tilde R_{ab} + \tilde R_{aZ}\,\varepsilon^{-1}v\,\tilde R_{Zb} , \qquad
  \varepsilon = 1 - vP ,
  \label{eq:D.27-screened}
\end{equation}
which is \cref{eq:8.4-screened} with $\varepsilon_{\vect{G}\vect{G}'} = \delta_{\vect{G}\vect{G}'} - v_{\vect{G}}\tilde R_{Z_{\vect{G}}Z_{\vect{G}'}}$. The same steps give the Dyson form $R_{ab} = \tilde R_{ab} + \tilde R_{aZ}\,v\,R_{Zb}$, since $R_{Zb} = (1 - Pv)^{-1}\tilde R_{Zb}$ by \cref{eq:D.27-charge}. For the electron density alone, $a = b = e$ and $n_Z = -n_e$, this is \cref{eq:4.4-dyson} and \cref{eq:4.4-epsGG}, and the one-species algebra of \cref{app:D.7} is the special case. The matrix $\varepsilon$ is invertible at every real $\omega$ except at the self-sustained longitudinal oscillations of the target, plasmons and longitudinal optical phonons, where $R$ itself has its poles.

With a single $\vect{G}$, $P$ and $v$ are numbers, and \cref{eq:D.27-screened} is \cref{eq:8.4-rank1}. Written as a matrix over the rows $a$, $b$ of the species, the screening term $v\tilde R_{aZ}\tilde R_{Zb}/\varepsilon$ is a rank-one correction of $\tilde R$, and \cref{eq:8.4-rank1} is the Sherman--Morrison form of the inverse in species space.

\emph{Total charge.} Take $a = b = Z$, so that $\tilde R_{aZ} = \tilde R_{Zb} = P$. Since $vP = 1 - \varepsilon$ and $1 - \varepsilon$ commutes with $\varepsilon^{-1}$,
\begin{equation}
  vR_{ZZ} = vP + vP\,\varepsilon^{-1}\,vP = (1 - \varepsilon)\bigl[1 + \varepsilon^{-1}(1 - \varepsilon)\bigr] = (1 - \varepsilon)\,\varepsilon^{-1} = \varepsilon^{-1} - 1 .
  \label{eq:D.27-ZZ}
\end{equation}
Its element at $\vect{G} = \vect{G}' = \vect{G}_q$ is $v_{\mathrm{C}}(q)R_{ZZ}(\vect{q},\omega) = (\varepsilon^{-1})_{00}(\vect{q},\omega) - 1$. With the fluctuation--dissipation theorem \cref{eq:4.2-fdt}, $S_{ZZ} = -2\Imag R_{ZZ}/(1 - \ee^{-\beta\omega}) = (2q^2/e^2)\,L/(1 - \ee^{-\beta\omega})$, with $L = \Imag[-(\varepsilon^{-1})_{00}]$ of the total dielectric matrix. No approximation enters: the relation holds for moving nuclei, any temperature and any $\omega$.

\emph{Electron density.} The electron density is $n_e = -n_Z + \sum_\kappa Z_\kappa n_\kappa$, and the response is bilinear in its two rows. Hence
\begin{equation}
  R_{ee} = R_{ZZ} - \sum_\kappa Z_\kappa\bigl(R_{Z\kappa} + R_{\kappa Z}\bigr) + \sum_{\kappa\kappa'}Z_\kappa Z_{\kappa'}R_{\kappa\kappa'} .
  \label{eq:D.27-ee}
\end{equation}
With the nuclei clamped, each $n_\kappa$ is a c-number Bragg sum, its commutators vanish, and every term with a nuclear row drops out: $R_{ee} = R_{ZZ}$, and the electron density is screened by the loss function of the total dielectric matrix, which is then the electronic one. With moving nuclei the nuclear terms describe the charge that the nuclei carry and the electrons they drag; they are what separates the electron-channel function $\varepsilon_e^{-1} = 1 + v_{\mathrm{C}}R_{ee}$ of \cref{app:A.7} from $\varepsilon^{-1}$.

\subsection{Exactness of the loss-function form}

Let the coupling be $O_c = \sum_sc_sn_s(\vect{q})$, with $c_s = c_1^{(s)}$ of \cref{sec:8.1}, real, independent of $\vect{q}$ (nuclear form factors set to one), and $S_{cc}$ its structure factor. By \cref{eq:D.27-ZZ}, \cref{eq:8.4-exact} is the statement $S_{cc} = \lambda_Z^2S_{ZZ}$ at every $\vect{q}$ and every $\omega \neq 0$; at $\omega = 0$ the static Bragg terms of a clamped lattice differ, and only the charge itself, with $D = 0$ below, holds there too. The target is a crystal: the small-momentum step below uses its long-wavelength longitudinal acoustic phonon.

\emph{Sufficiency.} Suppose $O_c - \lambda_Zn_Z = D$ has no matrix elements between states of different energy. In the Lehmann form \cref{eq:4.2-Sab}, $\langle f|O_c|i\rangle = \lambda_Z\langle f|n_Z|i\rangle$ for every pair with $E_f \ne E_i$, and $S_{cc} = \lambda_Z^2S_{ZZ}$ at every $\omega \ne 0$. For $c_e = -\lambda_Z$, $c_\kappa = \lambda_ZZ_\kappa$, $D = 0$. With clamped nuclei, $D = (c_e + \lambda_Z)n_e + \sum_\kappa(c_\kappa - \lambda_ZZ_\kappa)n_\kappa$ is static exactly when $c_e = -\lambda_Z$, for any $c_\kappa$; the sign of $\lambda_Z$ is free.

\emph{Necessity, large momentum.} Let the nuclei move. For $q$ large compared with the inverse size of every electron orbital and with $q^2/2M_\kappa$ large compared with the phonon energies, the impulse regime of \cref{sec:8.5}, every particle recoils as if free during the time $1/\omega$ that the response probes. The self term of each particle in $\langle O_c(\vect{q},t)^\dagger O_c(\vect{q},0)\rangle$ then gives a recoil peak centered at $q^2/2m_s$. Its width is of order $q\,p_s/m_s$, with $p_s$ the root-mean-square momentum of the particle in the target, and its weight is $\int\dd\omega\,S_{cc}/2\pi = c_s^2\nd_s$ per species, by the zeroth moment of the self term. The terms between two different particles involve $\langle\ee^{-\ii\vect{q}\cdot\vect{r}_i(t)}\ee^{\ii\vect{q}\cdot\vect{r}_j(0)}\rangle$, the Fourier transform of a pair correlation that is smooth on the scale $1/q$, and vanish as $q \to \infty$; in particular the electrons no longer follow the nuclei, and the electron--nucleus interference disappears. The distance between the peaks of two masses grows as $q^2$ and their widths as $q$, so at large enough $q$ each mass has a separate peak. Group the nuclear species by mass. Since nuclei of different atomic number have different masses, the nuclei of one group share an atomic number $Z_g$; let $\nd_g$ be their total density. Equality of $S_{cc}$ and $\lambda_Z^2S_{ZZ}$ peak by peak requires
\begin{equation}
  c_e^2 = \lambda_Z^2 , \qquad \sum_{\kappa\in g}c_\kappa^2\nd_\kappa = \lambda_Z^2Z_g^2\,\nd_g \quad\text{for every mass group } g .
  \label{eq:D.27-peaks}
\end{equation}

\emph{Necessity, small momentum.} At $\vect{q} \to 0$ and the frequency of a longitudinal acoustic phonon, the target is displaced locally as a whole: a long-wavelength acoustic phonon is a slow, local translation of the crystal, and the electrons go with it. A displacement field $\vect{U}(\vect{r})$ that varies on the scale $1/q$ moves each particle from $\vect{r}_j$ to $\vect{r}_j + \vect{U}(\vect{r}_j)$, and $n_s(\vect{q}) \to \sum_{j\in s}\ee^{\ii\vect{q}\cdot\vect{r}_j}[1 + \ii\vect{q}\cdot\vect{U}(\vect{r}_j)]$. For the mode of wavevector $\vect{q}$, $\vect{U}(\vect{r}) \propto \ee^{-\ii\vect{q}\cdot\vect{r}}$, and the amplitude to create the phonon is proportional to $q\sum_sc_sN_s$, with $N_s$ the number of particles of species $s$, up to corrections of relative order $q$. The weight of the acoustic peak in $S_{cc}$ is therefore $q^2(\sum_sc_s\nd_s)^2$ times a factor common to all couplings. For $n_Z$ the sum is the net charge-number density, which vanishes by neutrality, $\nd_e = \sum_gZ_g\nd_g$, and the acoustic weight of $S_{ZZ}$ starts at order $q^4$, from the charge that a long-wavelength strain induces. Equality as $q \to 0$ requires
\begin{equation}
  c_e\nd_e + \sum_\kappa c_\kappa\nd_\kappa = 0 .
  \label{eq:D.27-acoustic}
\end{equation}

\emph{Conclusion.} Choose the sign of $\lambda_Z$ so that $c_e = -\lambda_Z$ with $\lambda_Z > 0$; \cref{eq:D.27-acoustic} becomes $\sum_\kappa c_\kappa\nd_\kappa = \lambda_Z\nd_e = \lambda_Z\sum_gZ_g\nd_g$. By the Cauchy--Schwarz inequality with the weights $\nd_\kappa$, and by \cref{eq:D.27-peaks},
\begin{equation*}
  \sum_{\kappa\in g}c_\kappa\nd_\kappa \le \Bigl(\sum_{\kappa\in g}c_\kappa^2\nd_\kappa\Bigr)^{1/2}\nd_g^{1/2} = \lambda_ZZ_g\,\nd_g ,
\end{equation*}
with equality only if $c_\kappa = \lambda_ZZ_g$ for every species of the group. Summed over the groups, the left-hand side reaches $\lambda_Z\sum_gZ_g\nd_g$ only if equality holds in every group. Hence $c_e = -\lambda_Z$ and $c_\kappa = \lambda_ZZ_\kappa$ for every $\kappa$: $O_c = \lambda_Zn_Z$. With the nuclei clamped there are no nuclear recoil peaks and no acoustic phonon, and \cref{eq:D.27-peaks} leaves only $c_e^2 = \lambda_Z^2$, which sufficiency already covers. This proves \cref{eq:8.4-exact}.

The hypothesis on the masses excludes isobars of different atomic number in one target, which no crystal of this book contains. Without it, two species of equal mass and different $Z$ would share a recoil peak, and \cref{eq:D.27-peaks} would constrain only a combination of their couplings.

\subsection{Coulomb-line classes}

The classification of \cref{sec:8.4-lines} is a statement about the diagrams of the response at the random-phase level, also called time-dependent Hartree, at order $g^2$ in the rate with one phonon line. In this approximation every induced density acts on the electrons only through its Hartree field. A diagram is then built from fermion loops, each a Kohn--Sham response of some order. A loop takes its input fields, external, Hartree or ionic, and returns one induced density. That density reaches the rest of the diagram through one Coulomb line, or it is the output of the response. Each loop therefore has exactly one line toward the output, and the loops and Coulomb lines form a tree. Lines inside a loop, the exchange and ladder diagrams of \cref{sec:4.8}, are absent at this level.

Now add the phonon line, which joins two ion vertices on two loops of the diagram. Two cases arise.

\emph{The diagram without the phonon line is connected.} Then it is a single tree, and adding one line to a tree closes exactly one cycle. Every Coulomb line off the cycle is a bridge: cutting it separates a piece that is attached to the rest only through that line. Momentum conservation fixes its momentum to the total momentum entering that piece. The piece contains no part of the cycle and at most the two external vertices, so the momentum is $\vect{q}$ up to a reciprocal-lattice vector and the frequency is $\omega$. These are the external lines, and their chains are resummed by \cref{eq:D.27-screened}. On the cycle, let $A$ and $B$ be the loops at which the paths from the two external vertices meet it; a loop that carries an external vertex and lies on the cycle is its own meeting point. The cycle consists of the phonon line and a path from one ion vertex to the other that passes through $A$ and $B$. Momentum is conserved at every loop, and the only momenta that enter the cycle from outside are $\vect{q}$ at $A$ and $-\vect{q}$ at $B$ (or the reverse). Follow the cycle from the ion vertex at one end, with $\vect{Q}$ the momentum of the phonon line. The momentum is $\vect{Q} + \vect{G}$ up to the first of $A$ and $B$, it is $\vect{Q} + \vect{q} + \vect{G}$ between $A$ and $B$, and it is $\vect{Q} + \vect{G}$ again from the second of them to the other ion vertex. The frequencies follow in the same way: the frequency $\omega'$ of the phonon line on the segments next to the ion vertices, $\omega - \omega'$ or $\omega + \omega'$ between $A$ and $B$, and $\omega' = \pm\omega_{\vect{Q}\nu}$ when the phonon is on shell in the final state. The segments next to the ion vertices are the vertex lines, and the segment between $A$ and $B$ is the pair segment; when $A = B$, it is empty. Every Coulomb line is thus external, vertex or pair, and only one of these.

The chains on the vertex segments, summed over all diagrams, dress the bare potential of the displaced ion with every Coulomb chain at $(\vect{Q} + \vect{G}, \omega')$ before it reaches the first loop that carries an external vertex. This is the screened potential $\varepsilon^{-1}(\vect{Q},\omega')$ acting on the bare ionic potential, \cref{eq:6.2-screen} at the frequency $\omega'$. The chains on the pair segment between two loops sum to $v + vRv = \varepsilon^{-1}v$ at $(\vect{q} + \vect{Q}, \omega \mp \omega_{\vect{Q}\nu})$, the screened interaction $W^{\mathrm{scr}}$ of \cref{sec:4.8}.

\emph{The diagram without the phonon line falls into two trees.} Then the phonon line is a bridge, each tree contains one external vertex, and the phonon line carries $(\vect{q} + \vect{G}, \omega)$. It is not a Coulomb line but the ionic polarization inserted into the external chain, and summing such insertions turns the electronic dielectric matrix into the total one. This case does not describe a pair with a phonon in the final state: its phonon line is virtual and screens the dark-matter vertex. In the amplitude of the phonon-assisted process it would be a second phonon line, at $(\vect{q} + \vect{G}, \omega)$, beside the real phonon at $\vect{Q}$, and it lies outside the one-phonon diagrams of the first case. Using the total $\varepsilon$ of \cref{eq:D.27-screened} in the external and pair chains includes it once.

The assignment uses only the position of a line relative to the cycle and the external vertices, not the values of the momenta. It remains unique when the momenta coincide modulo $\vect{G}$, at $\vect{q} \to 0$, where the pair and vertex segments both carry $\vect{Q}$, or at $\vect{Q} \to 0$, where the pair segment and the external lines both carry $\vect{q}$.

\emph{The static vertex.} Density-functional perturbation theory evaluates the vertex chains at $\omega' = 0$. For densities the reality of the response in real space gives $R_{\vect{G}\vect{G}'}(\vect{Q},\omega')^* = R_{-\vect{G},-\vect{G}'}(-\vect{Q},-\omega')$, and time reversal of the equilibrium state gives $R_{\vect{G}\vect{G}'}(\vect{Q},\omega') = R_{-\vect{G}',-\vect{G}}(-\vect{Q},\omega')$. Together, $R(\vect{Q},-\omega') = R(\vect{Q},\omega')^\dagger$. Below the gap the response has no absorptive part, its anti-Hermitian part vanishes, and $R(\vect{Q},-\omega') = R(\vect{Q},\omega')$; the same holds for the proper response and for $\varepsilon$. The function is analytic for $|\omega'| < \Eg$, since its poles lie at the excitation energies, and an even analytic function differs from its value at zero by a term of order $\omega'^2$. With the scale $\Eg$ set by the nearest poles, the static vertex errs by a relative amount of order $(\omega_{\vect{Q}\nu}/\Eg)^2$. Without time reversal, in a magnetic insulator, only the reality relation remains, which ties $\omega'$ to $-\omega'$ at $-\vect{Q}$; the linear term survives, and the error is of order $\omega_{\vect{Q}\nu}/\Eg$.
